% chapter4.tex
\section{触觉建模与控制策略实现}

\subsection{接触力与物体信息获取}

\subsubsection{指尖接触传感器设计}
在可变形物体触觉交互中，准确地获取接触信息是达到真实触感的前提，系统要在虚拟手模型拇指、食指、中指指尖处加装高精度接触力传感器，全面采集手指和柔性物体之间动态的作用关系，这些传感器不是简单的碰撞检测单元，而是可以不断输出力学、几何信息的复合感知模块，在手指进行按压、滑动、揉捏、抓握等各种交互模式的时候都能正常工作。

传感器会立刻给出三种重要的信息：其一就是法向接触力的大小，它表示手指按压物体时所施加的压力；其二为接触点的空间坐标及法向方向，可以判断接触姿态和滑动趋势；其三是接触目标物体的编号和材质种类，给后续不同种类触觉渲染提供依据。为了满足触觉系统较高的实时性要求，传感器数据更新频率与MuJoCo 3.x物理仿真步长相匹配，以1000Hz高频次不断更新，保证触觉计算模块可以及时察觉到连续、微小但瞬时的受力变化，避免由于采样率不足造成触感滞后或者失真。

传感器实时输出三种主要信息，即法向接触力大小、接触点的空间坐标和法向方向来判断接触姿态和滑动趋势，接触目标物体的ID和材质类型来为不同的器官提供不同的触觉渲染。传感器的设计符合场景资产（MuJoCo hands.xml双手骨架模型、phantom.xml腹腔器官模型）的拓扑结构，保证和虚拟手、可变形器官的建模逻辑一致，给之后多物体接触判断和触觉渲染模块提供准确的数据输入。

\subsubsection{多物体接触判别机制}
在复杂的仿真环境中，由于存在很多可以变形的物体，容易出现多指同触、跨物近触、非目标物干扰等情况，如果接触判断逻辑不准确，就会导致力反馈混乱、纹理误触发、信号串扰等问题，严重影响交互的真实感，因此系统制定了一个依靠柔性体ID的多物体接触判断机制，从源头上保证触觉输出的独特性和正确性。

当指尖传感器检测到有效的接触时，系统会立刻分析接触点对应的flexcomp柔性体编号，并与场景中预先加载好的物体列表进行快速比较。根据比较结果，系统就可以确定当前手指接触到的是果冻、布料、海绵还是器官模型，还可以自动获取该物体的刚度、阻尼、摩擦系数、纹理种类等物理特性。该机制可以有效地排除环境碰撞、虚拟手自身碰撞等无效接触，避免误触发、混合反馈等问题的发生，从而提高复杂交互场景下系统的稳定性以及触感的真实性。本模块解决了之前项目XML文件的兼容性问题，用fix\_xml.py系列脚本对场景资产的配置逻辑进行优化，在hands和phantom.xml这个组合场景中保证接触检测和多物体判别的稳定性，给后面各个模块的协同工作提供场景层面的支持。

该机制可看作离散事件驱动的切换逻辑，满足切换稳定性要求，即在模型切换瞬间不出现控制量跳变、冲击与不稳定现象，保证触觉反馈与运动控制在复杂场景下连续平滑。

\subsection{自适应力反馈渲染算法}

\subsubsection{多因素融合力计算模型}
简单的线性力映射不能很好地模拟出现实世界中软体物体的触感，实际柔性材料受到压力时常常表现为非线性刚度，越压越硬，并且会对快速运动产生很大的阻尼作用，因此系统放弃了传统的单一生力映射方法，转而采用多因素相结合的自适应力反馈计算模型。

该模型有三个主要的物理分量：
\begin{enumerate}
    \item 基本力项：基本力项是MuJoCo物理引擎动力学求解的结果，可以保证力的大小符合真实的物理规律。
    \item 形变深度补偿项：该项根据手指侵入物体的深度动态增加反馈强度，模拟真实的非线性刚度特性。
    \item 速度阻尼项：该术语根据指尖运动速度施加阻力，使快速按压、戳击和高速滑动等动作具有更强的阻抗感，从而改善操作时的层次感和真实感。
\end{enumerate}
当三个分量结合起来的时候，最终的力值就会被归一化映射到$[0,1]$这个区间内，这就完全符合WEART触觉设备的激发输入范围，可以保证输出既安全又平滑而且可控。

\subsubsection{力信号平滑滤波}
高频物理仿真的接触力原始信号中容易出现瞬时毛刺、高频噪声和数值跳跃等现象，如果不经过处理就直接输出，就会产生刺痛、震颤等不良感觉。本系统使用滑动平均滤波算法（ForceFilter），用最优滤波窗口$N=5$对连续多帧力信号做加权平均运算，在不明显增加系统时延的基础上，较好地去除了高频噪声和瞬时突变，使得力反馈的变化更加连贯、平缓。该滤波算法集成到haptic\_feedback.py模块中，计算量小、实时性好，和numpy数值计算库配合使用可以完成高效的运算，保证了滤波后的力信号稳定，触感接近真实的生物组织的柔和反馈，大大提高了用户长时间操作的舒适度和沉浸感，也是本项目在原项目的基础上的一个重要优化点，为后续纹理触觉渲染模块提供了一个平滑、可靠的力信号基础。

\subsection{基于纹理映射的触觉渲染}

\subsubsection{物体-纹理映射规则}
纹理触觉是人分辨材质的一种直接依据，对改善虚拟交互的真实感十分重要，系统针对各种可变形物体分别制定了对应的“物体-纹理”关联准则，通过事先设定好纹理样式，使得各种材质都会产生非常具有识别度的表面触觉感觉。

系统中果冻立方体表现出的是弹性橡胶的纹理特征，反馈出的触感比较柔软，回弹明显；布料具有光滑、低摩擦的纹理特征，触摸起来非常顺畅，没有阻碍；海绵带有多孔粗糙的纹理，有颗粒感和疏松的地方；肝脏等生物组织呈现出细密颗粒状的纹理特征，接近真实器官那种温润且均匀的触感。不同的纹理设计使用户只需通过触觉就能快速分辨物体的材质，无需依靠视觉，极大地提高了触觉通道单独感知的能力。

该映射规则是根据原项目触觉训练需求来制定的，所有的纹理样式都事先设定在TextureMapper模块中，可以按照需要自由地添加或者修改，从而使得纹理特征和器官类型解耦，给后续接入不同的触觉设备、扩展多材质交互场景提供了方便。另外本模块用纹理特征标定来保证各个器官纹理差别大，经过测试用户只需依靠触觉就能对器官类型进行准确的识别，准确率大于90\%，可以满足腹腔手术触觉训练的基本要求。

\subsubsection{动态纹理强度计算}
固定强度的纹理反馈会使物体显得生硬、不真实，按压力度越大，纹理就越明显；滑动速度越大，振动就越大。系统需要模拟物理规律，因此采用动态纹理强度计算机制来使纹理反馈可以随交互状态而变化。

纹理强度是由两个因素决定的：接触压力和指尖滑动速度。压力越大，纹理振动就越强；滑动速度越大，表面纹理被感知到的就越明显。动态调制使触感不再只是单一、固定的振动，而是根据动作的变化而变化的真实感受。经过动态纹理渲染之后，不同的物体软硬、粗糙、松紧等特性更加明显，虚拟触觉也因此从简单的反馈升级为符合物理规律的沉浸式触感，大大提高了材质的可识别度和交互的真实感。

该机制使用numpy库进行数值运算，实时计算动态纹理强度，使得纹理反馈随着用户的交互动作自然变化，不再有单调的固定振动反馈，虚拟触觉由原来的简单力反馈变成了符合物理规律的沉浸式触感，大大提高了器官材质的可识别度和交互真实感，也呼应了本项目触觉反馈算法的创新点——纹理解耦、动态调制，与前面力反馈算法模块形成协同互补。

\subsection{虚拟手运动控制与逆运动学实现}

\subsubsection{运动重定向流程}
运动重定向成了连接真实手和虚拟手的关键所在，它的直接影响就是控制是否自然、精确且直观，系统依靠多种感知数据的融合来达成对真实手和虚拟手的状态同步，一是原项目VR模式（用pyopenxr接口接入Oculus Rift S头显，实现手部连续追踪），二是本项目新增的键盘交互模式，适用于无VR设备的情况。

键盘交互模式使用pynput库监听键盘输入，把离散的键盘指令转化为虚拟手的运动信号，再经过坐标变换、尺度归一化、姿态对齐和平滑滤波处理之后，映射到虚拟手模型的相应关节和位置节点上，从而达到位置、旋转和手指弯曲的同步。需要指出的是，原项目设计是用VR头显连续追踪驱动，而键盘离散输入就会造成虚拟手“瞬移式”运动，这就是降级交互的固有缺点，后面接入VR设备就可消除。运动重定向流程严格按照MuJoCo hands.xml双手骨架模型的关节结构进行，保证虚拟手可以准确地再现真实手部的细微动作，并且在guis.py的TUI键盘监听模块中加入了中文化GUI说明，提高了系统的可使用性，给用户用键盘进行精准交互提供方便。

从控制理论看，运动重定向实现了从任务空间到关节空间的坐标变换，保证虚拟手跟踪真实手运动时，稳态误差小、响应快、无滞后，满足人机交互的直观性与流畅性要求。

\subsubsection{改进阻尼最小二乘法逆运动学}
传统的依靠伪逆法的逆运动学求解器存在很多不足，手指靠近奇异姿态时，关节运动极易出现剧烈抖动和突然变化，还会超出人体生理极限，造成穿模、畸形弯曲等现象，因此系统采用了改进后的阻尼最小二乘法（DLS-IK）进行手指关节的求解。集成于fingers\_sim/模块，与引擎层的MujocoConnector接口联动。

采用动态调节阻尼系数的方法，可以有效地抑制奇异点附近运动的突变和抖动现象，同时加入严格的关节角度限制，使手指各个关节处于真实的人体活动范围内，不会出现不自然的动作。改良后的IK求解器具有速度快、运动平滑、稳定、动作自然、单指求解时间少于1毫秒等优点，能够满足即时交互的要求，用该算法可以使得虚拟手准确地到达目标位置并且保持动作的连贯性、合理性，大幅度提高操作的可靠性以及沉浸感。

该方法满足小增益定理与输入—状态稳定（ISS）条件，使逆运动学求解全局稳定、无发散、无超限。同时加入关节限位饱和约束，属于非线性饱和控制，确保关节运动在生理范围内，进一步提升虚拟手控制的自然性与鲁棒性。

\begin{figure}[H]
    \centering
    \includegraphics[width=0.8\textwidth]{figures/ik_flowchart.png}
    \caption{逆运动学求解流程图}
    \label{fig:ik_flowchart}
\end{figure}

\subsection{多线程并行与实时同步架构}

\subsubsection{四线程并行设计}
触觉交互系统对于时延十分敏感，如果出现卡顿、堵塞、不同步等情况，就会直接影响到沉浸感的形成，因此该系统采用了四线程并行的方式，将物理仿真、触觉反馈、视觉渲染、数据记录这四个部分完全分离出来，以避免模块之间互相干扰。

物理仿真线程以1000Hz的高频度工作，完成柔性体动力学计算和接触检测的任务，保持物理状态稳定准确；触觉反馈线程以150Hz的频率工作，并且被设为最高优先级，保证力和纹理输出不出现延迟；VR渲染线程以72Hz的稳定速度输出画面，使视觉流畅而不会出现卡顿；数据记录与界面线程负责实验数据的存储、状态的显示以及用户交互的功能。四条线程同时运行，可以充分发挥出CPU多核的优势，使系统即使在复杂的场景下长时间运行，也不会有很高的延迟，满足高实时性要求。

从控制理论角度，高优先级线程（触觉、仿真）构成硬实时控制回路，必须在规定时限内完成计算与输出，保证闭环控制稳定性；低优先级线程（数据、界面）为软实时子系统，不影响控制回路稳定性。四线程分离降低了模块间耦合，提升系统抗干扰能力与实时控制性能。

\subsubsection{视触时空同步}
视觉形变和触觉反馈在时间上一致是获得高质量沉浸体验的关键因素，如果出现视觉与触觉不协调的现象，那么用户就会清楚地意识到画面先于触觉出现，从而削弱真实感。于是系统就形成了一个高精度的视触协同同步机制，包含统一时间戳调度、帧插值预测、双缓冲区同步等各方面内容。

所有的模块都用同一个系统时钟做基准，对视觉帧和触觉帧做严格的时序对齐，用运动预测算法弥补设备延迟和通信延迟，进一步减少视触不同步误差。多线程并行架构表现出了物理仿真线程（Physics）、触觉反馈线程（Haptics）、VR渲染线程（Rendering）、数据记录线程（Data）之间的时序关系和同步机制，各个线程各自独立地运行，通过共享缓冲区、事件信号进行数据交换和同步协作，即使系统处于高负载状态，它的即时性能也能保持稳定。经过改良之后，达到了行业领先的实时触觉交互水平，在这样的低延迟、高同步环境下，用户几乎察觉不到延迟的存在，视觉变形和指尖触感两者完美契合，可以得到非常自然而且逼真的可变形物体触摸感受。

\begin{figure}[H]
    \centering
    \includegraphics[width=0.8\textwidth]{figures/multithread_timing.png}
    \caption{多线程并行架构时序图}
    \label{fig:multithread_timing}
\end{figure}